\documentclass[10pt,a4paper]{amsart}
\usepackage[margin=19mm]{geometry}
\usepackage{amsmath,amssymb,mathtools}
\usepackage[scr=rsfs,cal=euler]{mathalfa}
\usepackage{xfrac}
\usepackage[unicode,hidelinks,bookmarks=false]{hyperref}
\newcommand{\ie}{i.e.\ }
\newcommand{\cf}{cf.\ }
\newcommand{\resp}{respectively\ }
\newcommand{\Subsection}{\subsection}
\providecommand{\nobreakdash}{\nobreak\hbox{-}}
\setlength{\emergencystretch}{2em}
\multlinegap0pt
\usepackage{bm}
\usepackage{bbm}
\usepackage{dsfont}
\usepackage{xspace}
\usepackage{cancel}
\usepackage{mathtools}
\usepackage{amsthm}
\makeatletter
\@ifundefined{theorem}{\newtheorem{theorem}{Theorem}}{}
\@ifundefined{lemma}{\newtheorem{lemma}[theorem]{Lemma}}{}
\makeatother
\DeclareMathOperator{\ord}{ord}
\def\Q{\mb Q}
\def\Z{\mb Z}
\def\k{\mb K}
\def\mb{\mathbb}
\def\wt{\widetilde}
\title[On Goncharov's conjecture in next to Milnor degree]{On Goncharov's conjecture in next to Milnor degree}
\author{Vasily Bolbachan}
\date{Source: arXiv:2404.06271v2 (CC BY 4.0)}
\begin{document}
\maketitle

\noindent\emph{Source and licence.} On Goncharov's conjecture in next to Milnor degree. Version arXiv:2404.06271v2 (CC BY 4.0), distributed under the Creative Commons Attribution 4.0 International licence, as recorded in the arXiv licence metadata for this exact version and shown on the abstract page https://arxiv.org/abs/2404.06271v2. Full rights notice: licenses/arxiv-2404.06271-CC-BY-4.0.txt.\par\smallskip

\noindent\textit{Editorial scope.} Counted unit: the Lemma whose source label is \texttt{lemma:linear\_independence} in the article (source0.tex lines 1388--1392, source label \texttt{lemma:linear\_independence}), delivered together with the article's own complete original proof of it (source0.tex lines 1394--1430, one \texttt{proof} environment of 37 lines). No mathematics is written, repaired or re-derived: statement and proof are the article's own text line for line. Required context retained uncounted: prop:theta\_tensor\_is\_injective (lines 1357--1360); lemma:about\_basis (lines 1374--1378). Every definition, display and label of those lines is retained unchanged. Omitted from this package and not redistributed: 1--1356, 1361--1373, 1379--1387, 1393--1393, 1431--2106. Nothing inside a retained excerpt is omitted or altered: every retained body is delivered exactly as the article's lines are written, so no omission receipt of any kind accompanies this package. Citations: the retained excerpts carry no citation command at all, so no bibliography entry of the article is transcribed and no reference is redistributed. Reference adaptation: none was needed and none was made; every reference inside the delivered unit resolves inside it, so no unresolved reference or citation is produced. Printed slips kept literal: no printed word, symbol, hypothesis or number of the counted statement or of its proof was altered. Layout and class: the article's own class is \texttt{extarticle} with packages including amsmath, amssymb, amsthm, geometry, comment, tikz-cd, csquotes; the wrapper is the lane's pinned \texttt{amsart} editorial wrapper at 10pt on A4, which restates the theorem-like environments the retained text needs and adds the article's own macro definitions that the retained text needs (\texttt{\textbackslash Q}, \texttt{\textbackslash Z}, \texttt{\textbackslash k}, \texttt{\textbackslash mb}, \texttt{\textbackslash ord}, \texttt{\textbackslash wt}), with extra wrapper packages (bm, bbm, dsfont, xspace, cancel, mathtools). The delivered numbering is the unit's own. Assets: no retained span contains a figure environment, a graphics-inclusion command, a PDF-page inclusion, a table or a TikZ picture; no image file is copied into this package, the package declares no asset dependency and no image file is redistributed with it. Licence: version arXiv:2404.06271v2 of this article is distributed under the Creative Commons Attribution 4.0 International licence, as recorded in the arXiv licence metadata for this exact version and shown on the abstract page https://arxiv.org/abs/2404.06271v2. A full rights notice is delivered at licenses/arxiv-2404.06271-CC-BY-4.0.txt. Characters: every retained excerpt is pure ASCII, so the delivered engine, XeLaTeX with its default Latin Modern text font, reports no missing character.
\begin{lemma}
\label{prop:theta_tensor_is_injective}
    The map $\theta_{\otimes,Z,n}$ is injective.
\end{lemma}

\begin{lemma}
\label{lemma:about_basis}
    Let $X$ be a smooth projective curve and $Z\subset X$ be a finite subset of $X$. There is a basis $h_\beta$ of the vector space $\k(X)^\times\otimes_\Z\Q$ with the following property.  Assume that for some $f\in \k(X)^\times$ we have the decomposition $$f=\prod_\beta  h_\beta^{n_\beta},\quad n_\beta\in \Q.$$
Assume also that for some $z\in Z$, we have $\ord_z(f)=0$. Then for any $\beta$ such that $n_\beta\ne 0$ we have $\ord_z(h_\beta)=0$.
\end{lemma}

\begin{lemma}
\label{lemma:linear_independence}
    Let $V,W$ be vector spaces. Let $v_1,\dots, v_k\in V, w_1,\dots, w_k\in W$. Assume that $v_i$ are linearly independent. If
    $\sum v_k\otimes w_k=0$ then $w_k=0$.
\end{lemma}

\begin{proof}[The proof of proposition \ref{prop:theta_tensor_is_injective}]
    The proof is by induction on $n$. The case $n=1$ is obvious.

    Let $x\in T_{ Z,n}(X)$ such that $\theta_{\otimes,Z,n}(x)=0$. We need to check that $x=0$. Let
    $$x=\sum_{\alpha\in A} n_\alpha\langle f_1^\alpha\wedge\dots\wedge f_n^\alpha\rangle_\otimes.$$
    We can assume that $n_\alpha\ne 0$.

    Let $Q$ be the set $$\sup\left(\sum_{i,\alpha}\sup((f_i^\alpha))\right)$$ and let $\wt Z=Z\cup Q$. Let $h_\beta$ be a basis from Lemma \ref{lemma:about_basis} for $Z=\wt Z$.

    Replacing $x$ with $Nx, N\in\mb Z$ we can assume that coefficients of each $f_1^\alpha$ in this basis are integers. So for some $l_{\alpha,\beta}\in\Z$ we have
    $$f_1^\alpha=\prod {h_\beta}^{l_{\alpha,\beta}}.$$

    For fixed $\alpha$ denote by $B_\alpha$ the set of $\beta$ such that $l_{\alpha,\beta}\ne 0$. For fixed $\beta$ denote by $A_\beta$ the set of $\alpha$ such that $l_{\alpha,\beta}\ne 0$. Let $B=\cup_\alpha B_\alpha$. We get the following relations in the group $T_{Z,n}(X)$:

    $$\langle f_1^\alpha,\dots, f_n^\alpha\rangle_\otimes=\sum_{\beta\in B_\alpha} l_{\alpha,\beta}\langle h_\beta,f_2^\alpha,\dots, f_n^\alpha\rangle_\otimes.$$

    So we get:

    $$x=\sum_{\alpha\in A} n_\alpha\sum_{\beta\in B_\alpha} l_{\alpha,\beta}\langle h_\beta,f_2^\alpha,\dots,f_n^\alpha\rangle_\otimes=$$
    $$\sum_{\beta\in B} \sum_{\alpha\in A_\beta} n_\alpha l_{\alpha,\beta}\langle h_\beta,f_2^\alpha,\dots,f_n^\alpha\rangle_\otimes.$$

    Let $Z_\beta$ be the union of $Z$ and $\sup((h_\beta))$. Denote by $y_\beta$ the following element in the group $T_{ Z_\beta, n-1}$:
    $$y_\beta = \sum_{\alpha\in A_\beta} n_\alpha l_{\alpha,\beta}\langle f_2^\alpha,\dots,f_n^\alpha\rangle_\otimes.$$

We get:

    $$\theta_{\otimes, Z,n}(x)=\sum_{\beta\in B} h_\beta\otimes (\theta_{\otimes, Z_\beta, n-1}(y_\beta)).$$

    By previous lemma, we get $\theta_{\otimes, Z_\beta, n-1}(y_\beta)=0$. It follows from the inductive assumption that $y_\beta=0$.

Assume that $\beta\in B$. For any $z\in Z$ we have $\ord_z(h)=0$. Define the map $L_\beta\colon T_{Z_\beta, n-1}(X)\to T_{ Z, n}(X)$ given by the formula
$$\langle g_2,\dots, g_n\rangle_\otimes\mapsto \langle h_\beta, g_2,\dots, g_n\rangle_\otimes.$$

We get

$$x=\sum_{\beta\in B} L_\beta(y_\beta)=\sum_\beta L_\beta(0)=0.$$
\end{proof}

\end{document}
